% math-spine/reversible-narrative-operators.tex
% Shared-Horizon Preprint — Mathematical Spine
% Module 3: Reversible Narrative Operators

\section{Reversible Narrative Operators}
\label{sec:reversible-narrative-operators}

\subsection{Narrative State Space}

Let $\mathcal{N}$ denote the narrative state space associated with Reversible
Narrative Ecology (RNE). We model $\mathcal{N}$ as a smooth expressive manifold
equipped with a narrative metric $g_{\mathcal{N}}$:

\[
(\mathcal{N}, g_{\mathcal{N}}) \quad \text{is a reversible narrative manifold.}
\]

The metric $g_{\mathcal{N}}$ encodes narrative coherence, reversibility, and
expressive stability.

\subsection{Reversible Narrative Operators}

A reversible narrative operator is a smooth map
\[
R : \mathcal{N} \rightarrow \mathcal{N}
\]
such that:

\begin{enumerate}
    \item $R$ is smooth: $R \in C^\infty(\mathcal{N})$.
    \item $R$ is reversible: $R^{-1}$ exists and is smooth.
    \item $R$ preserves narrative curvature:
    \[
    K_{\mathcal{N}}(R(p)) = K_{\mathcal{N}}(p).
    \]
\end{enumerate}

These conditions distinguish RNE from symbolic ambiguity, which lacks
reversibility and curvature preservation.

\subsection{Narrative Transitions}

A narrative transition is a reversible expressive map
\[
\tau : \mathcal{N} \rightarrow \mathcal{N}
\]
satisfying:

\[
\tau^{-1} \text{ exists and is smooth}.
\]

Narrative transitions must preserve expressive stability:

\[
E(\tau(p)) \le E(p) \quad \forall p \in \mathcal{N}.
\]

This ensures that narrative dynamics remain stable under forward and reverse
transformations.

\subsection{Reversible Narrative Flow}

A reversible narrative flow is a one-parameter family of operators
\[
\Phi_t : \mathcal{N} \rightarrow \mathcal{N}, \quad t \in \mathbb{R}
\]
satisfying:

\begin{enumerate}
    \item $\Phi_0 = \text{id}$.
    \item $\Phi_{t+s} = \Phi_t \circ \Phi_s$.
    \item $\Phi_t$ is reversible for all $t$.
\end{enumerate}

The flow is generated by a narrative vector field $X$:

\[
\frac{d}{dt} \Phi_t(p) = X(\Phi_t(p)).
\]

This defines reversible narrative evolution over time.

\subsection{Narrative Geodesics}

A curve $\gamma : [0,1] \rightarrow \mathcal{N}$ is a narrative geodesic if:

\[
\nabla_{\dot{\gamma}} \dot{\gamma} = 0
\]

with respect to the narrative metric $g_{\mathcal{N}}$.

Geodesics represent minimal-change narrative transitions and are used to compare
narrative states without activating operator geometry.

\subsection{Narrative–Expressive Compatibility}

Narrative operators must be compatible with expressive transitions:

\[
R \circ T = T \circ R
\]
whenever $T$ is an expressive transition on $\mathcal{N}$.

This ensures that narrative reversibility does not conflict with expressive
curvature constraints.

\subsection{Narrative–Operator Separation}

To prevent collapse between narrative and operator geometries, we require:

\[
O \circ R \notin \text{Domain}(O)
\]
for any operator $O$ defined on $\mathcal{M}$.

This condition preserves lineage sovereignty and prevents conceptual collapse.

\subsection{Reversible Narrative Operator Summary}

Reversible narrative operators define the mathematical backbone of RNE. They
provide reversible, curvature-preserving transformations that distinguish RNE
from symbolic ambiguity. These operators ensure narrative stability, expressive
compatibility, and separation from operator geometry.
